\section{功能条件生成：Motif Scaffolding 与生物学提示}

\subsection{从无条件生成到功能控制}

EvoDiff 框架的真正威力在于\textbf{功能条件生成} (conditional generation)——不仅能从头生成蛋白质，还能根据指定的功能约束来引导生成过程。这得益于离散扩散框架的天然灵活性。

核心思想是 \textbf{Motif Scaffolding}（功能基序的支架设计）：

\begin{importantbox}{Motif Scaffolding 的核心概念}
许多蛋白质的关键功能由序列中一个\textbf{局部的子区域}（称为 motif 或功能基序）所决定。例如，一个钙离子结合蛋白中，只有特定的几个氨基酸残基直接参与与钙离子的物理结合。

\textbf{Motif Scaffolding 的目标}：给定一个已知的功能 motif，设计一个全新的蛋白质序列来"支撑"（scaffold）这个 motif，使其能够在正确的三维空间位置上发挥功能。

这类似于自然语言处理中的 \textbf{inpainting}（填充）任务：给定序列中的一部分内容，让模型补全其余部分。
\end{importantbox}

\subsection{Inpainting 作为功能提示}

EvoDiff 的 Motif Scaffolding 能力直接源于其训练目标。由于模型学习了\textbf{所有可能的 masking 模式}的去噪，它自然具备了以下能力：

\begin{enumerate}
    \item 接收一条\textbf{部分已知}的序列——已知部分是功能 motif，其余位置为 \texttt{[MASK]}
    \item 仅对 mask 位置进行逐步去噪，保持 motif 部分不变
    \item 最终生成一条包含原始 motif 但其余部分全新的蛋白质序列
\end{enumerate}

Amini 将这种能力类比为\textbf{"生物学提示"} (biological prompting)——功能 motif 就是给模型的"提示"，模型基于这个提示生成满足功能约束的完整蛋白质。

\subsection{案例研究：钙离子结合蛋白的设计}

Amini 详细展示了一个具体的功能设计案例：

\textbf{背景}：天然存在一种蛋白质，其功能是结合钙离子 (Ca$^{2+}$) 并传递细胞信号。该蛋白质的完整序列中，有一段特定的子序列（绿色标记的 motif）直接负责与钙离子的物理结合和相互作用。

\textbf{设计过程}：
\begin{enumerate}
    \item 提取钙离子结合 motif 的序列片段作为"提示"
    \item 将 motif 放置在目标序列的对应位置，其余位置设为 \texttt{[MASK]}
    \item 运行 EvoDiff 的条件生成流程，模型逐步填充 mask 位置
    \item 得到一条包含原始钙结合 motif 但整体序列全新的蛋白质设计
\end{enumerate}

\textbf{实验验证}：
\begin{itemize}
    \item 将 EvoDiff 设计的蛋白质（蓝色曲线）与天然蛋白质（灰色曲线）进行钙离子结合实验
    \item 结果表明，AI 设计的蛋白质\textbf{成功展现出钙离子结合能力}
    \item 这意味着 EvoDiff 不仅能生成结构可行的蛋白质，更能生成\textbf{功能可行的}蛋白质
\end{itemize}

\begin{warningbox}{功能验证仍处于早期阶段}
Amini 坦诚指出，当前的功能验证实验虽然令人兴奋，但仍然是初步的"概念验证" (proof of concept)。从单个案例的成功到系统性的功能设计能力，还需要大量的后续工作。蛋白质设计的实验验证周期长、成本高，每一个成功案例背后可能有多个失败的尝试。
\end{warningbox}

\subsection{本章小结}

EvoDiff 的离散扩散框架天然支持功能条件生成（Motif Scaffolding）——通过将已知的功能 motif 作为"生物学提示"，指导模型生成包含该 motif 的全新蛋白质。钙离子结合蛋白的设计案例表明，EvoDiff 生成的蛋白质不仅结构可行，更能实现目标功能。这标志着 AI 驱动的功能性蛋白质设计迈出了重要的第一步。


